% OWNER: media-lab
% STATUS: outline
% SOURCES: notes/SPIS_TRESCI_SZCZEGOLOWY.md §4; GLOSSARY.md §8.3
\chapter{Sztuczna inteligencja dyfuzyjna dla architektury wnętrz --- medium i praktyka}
\label{ch04:top}
\ChapterStatus{outline}{media-lab}

\begin{quote}
Eksperymenty w ComfyUI (w tym kontrola geometrii przez mapy głębokości, conditioning oraz segmentację SAM/SAM2) stanowią warstwę praktyki autora; produkcyjny artefakt IDA realizuje generację przez syntezę promptu i image-to-image w środowisku Vertex~AI / Gemini Image (wcześniej FLUX na Modal).
\end{quote}

\section{Generatywne modele obrazu jako medium wizualizacji wnętrz}
\label{ch04:sec:genai-medium}
\Todo{Survey / kluczowe papers 2022--2026 o diffusion w architekturze i designie --- tylko zweryfikowane DOI.}

\section{Od text-to-image do image-to-image: geometria pokoju}
\label{ch04:sec:t2i-i2i}
\Todo{Zachowanie layoutu, perspektyw, otworów; typowe failure modes.}

% \begin{figure}[htbp]
%   \centering
%   \includegraphics[width=0.9\textwidth]{comfyui/ch04-geometry-fail-success.png}
%   \caption{Porównanie udanej i nieudanej geometrii pomieszczenia po restylizacji (opracowanie własne autora --- pracownia generatywna).}
%   \label{fig:ch04:geometry}
% \end{figure}

\section{Pracownia ComfyUI}
\label{ch04:sec:comfyui}

\subsection{Dobór i porównanie modeli}
\label{ch04:sec:modele}
\Todo{Tabela kryteriów: wierność geometrii, styl, czas, kontrola (SDXL / FLUX / warianty edycyjne).}

\subsection{LoRA i style adapters}
\label{ch04:sec:lora}
\Todo{Kontrola charakteru wizualizacji vs wierność zdjęciu bazowemu.}

\subsection{Workflow img2img i multi-reference}
\label{ch04:sec:img2img}
\Todo{Grafy ComfyUI; inspiracje jako referencje.}

\section{Kontrola struktury przestrzeni}
\label{ch04:sec:struktura}

\subsection{Mapy głębokości (depth)}
\label{ch04:sec:depth}
\Todo{Znaczenie depth dla układu pomieszczenia; typowe błędy.}

\subsection{ControlNet i conditioningi (canny, normals)}
\label{ch04:sec:controlnet}
\Todo{Conditioning struktury obrazu w kontekście wnętrz.}

\subsection{Segmentacja SAM / SAM2}
\label{ch04:sec:sam}
\Todo{Maski mebli, ścian, stref; restylizacja selektywna.}
\NeedsCite{Papers SAM (Kirillov et al.) oraz SAM~2 --- zweryfikować DOI przed cytowaniem.}

\section{Inpainting, empty shell i usuwanie wyposażenia}
\label{ch04:sec:inpaint}
\Todo{Od maski w pracowni do produktowego empty shell w IDA (prompt systemowy, bez pełnego ControlNet w runtime).}

\section{VLM i opis przestrzeni jako brief projektowy}
\label{ch04:sec:vlm}
\Todo{Od zrozumienia zdjęcia do briefu; most do analizy pokoju/inspiracji w IDA (Gemini Flash Lite).}

\section{Krytyka medium}
\label{ch04:sec:krytyka}
\Todo{Halucynacje geometrii; estetyka AI look; autorstwo i odpowiedzialność projektanta.}

\section{Checklista jakości wizualizacji koncepcyjnej}
\label{ch04:sec:checklista}
\Todo{Autorskie kryteria: geometria, światło, materiał, zgodność z briefem, czytelność funkcji.}
